\section{Monetary policy and the adoption margin}
\label{sec:monetary}

Two questions connect monetary policy to the adoption margin, and they point in
opposite directions. Does the margin reshape how a monetary shock transmits? And
does the monetary stance, during an energy shock, reshape the adoption wave? The
first is a robustness check; the second is where the interest-rate sensitivity of
green durables that \citet{dietrich2025consumer} identify in a representative agent shows up
in our setting.

\subsection{Monetary transmission}
\label{sec:mp_transmission}

The energy-shock experiments hold monetary policy at the AMRS constant-real-rate
baseline. A $100$-basis-point (annualised) tightening isolates the role of the
rule itself. Figure~\ref{fig:monetary} contrasts the constant-real-rate baseline
with a standard Taylor rule $(\phi_\pi,\phi_{\pi e})=(1.5,0)$. The response is
standard: output and consumption fall, inflation falls. It is \emph{not}
materially reshaped by the adoption margin. At a $5\%$ green steady state the
durable-choice channel is too small to move aggregate monetary transmission (on
impact, output $-0.86$ under the Taylor rule against $-1.13$ under the
constant-real-rate rule; inflation $-0.43$ against $-0.50$ annualised). The
adoption margin matters for how the economy absorbs an \emph{energy} shock, not
for how it transmits a \emph{monetary} one.

\begin{figure}[H]\centering
\figtitle{Monetary-policy shock}
\IfFileExists{output/fig_monetary_import.pdf}{\includegraphics[width=\textwidth]{output/fig_monetary_import.pdf}}{\textit{[figure output/fig\_monetary\_import.pdf: run run\_summary\_table.py]}}
\label{fig:monetary}
\fignotes{$100$bp annualised tightening: Taylor rule (solid) against the
constant-real-rate baseline (dashed). Panels: nominal rate, output, inflation,
consumption.}
\end{figure}

\subsection{The monetary stance and the adoption wave}
\label{sec:mp_stance}

The reverse question is Dietrich's. In their representative-agent economy green
durables are interest-rate sensitive, so a strict inflation target slows the
transition. The mechanism survives here and acquires a distributional dimension: the
switch is a lumpy purchase financed against the borrowing constraint, so its user
cost rises with the real rate, and it does so most for the lower-wealth households
that carry the crisis-induced adoption margin (Figure~\ref{fig:switch_heatmap}).

The stance operates on the wave, but weakly. Holding
the brown-price shock fixed and varying the rule, a Taylor rule that leans against
the energy-driven inflation ($\phi_\pi=1.5$) raises the ex-ante real rate, by $+4.5$
cumulated over $H=24$, and shrinks the peak green-share response from $8.03$ to
$7.87$ pp. A deliberate $25$-basis-point accommodation does the opposite, lifting it
to $8.08$ pp and cutting the output loss; the standalone rate cut raises adoption on
its own ($D^G$ peak $+0.09$ pp, consumption $+5.8$ cumulated) as easier financing
relaxes the constraint on the switch. The sign is Dietrich's; the magnitude is
small.


Figure~\ref{fig:monetary_stance} sweeps the inflation coefficient. Stricter
targeting shrinks the wave monotonically, from $7.93$ pp at $\phi_\pi=1.25$ to
$7.73$ at $\phi_\pi=3$, but the whole range spans a fifth of a percentage point,
while the cumulative output loss grows from $-70$ to $-91$. Strict inflation
targeting during an energy crisis buys a marginal slowdown of the transition at a
first-order output cost.

\begin{figure}[H]\centering
\figtitle{The monetary stance and the adoption wave}
\IfFileExists{output/fig_monetary_stance_import.pdf}{\includegraphics[width=\textwidth]{output/fig_monetary_stance_import.pdf}}{\textit{[figure output/fig\_monetary\_stance\_import.pdf: run run\_monetary\_stance.py]}}
\label{fig:monetary_stance}
\fignotes{Brown-price shock, as the Taylor inflation coefficient $\phi_\pi$ rises
(with $\phi_{\pi e}=0$). Left: peak $\Delta D^G$. Right: cumulative output. Solid:
Taylor rule swept over $\phi_\pi$; dashed: constant-real-rate baseline.}
\end{figure}

The comparison that matters is with the fiscal instruments. The peak wave's
semi-elasticity to the real rate is about $-0.035$ pp per unit of cumulative rate,
so moving the wave by one percentage point would take a cumulative rate change of
order thirty, an implausible monetary action, whereas the retail price cap removes
almost the entire $8.0$ pp wave directly (Section~\ref{sec:cap}). Monetary policy
reaches the adoption margin only through the financing cost of the switch, a
second-order channel; the cap and the transfer reach it through the relative price
itself, which is first-order. This is why we keep monetary policy at the
constant-real-rate baseline in the main text: not because the stance is neutral for
adoption, but because it is dominated, by an order of magnitude, by the instruments
the paper ranks.

\paragraph{An accommodative offset.} A natural crisis exercise runs an accommodative
innovation \emph{alongside} the energy shock, leaning against it rather than waiting.
Figure~\ref{fig:monetary_offset} sizes the accommodation to halve the cumulative
output loss, which takes a large front-loaded cut (about $144$ basis points
annualised). It softens the fall in cumulative output ($-65$ to $-32$) and consumption
($-109$ to $-76$) and lifts the adoption wave ($8.0$ to $8.4$ pp), the financing
channel adding to the effect, but pays for the relief with more inflation (impact annualised
inflation $+5.6\%$ to $+8.4\%$) and an output path that overshoots on impact before
settling. The offset buys output stabilisation and a small transition dividend at the
cost of the price level and a much larger rate move than the fiscal instruments
require; it does not substitute for the relative-price instruments, which stabilise
output without that inflation cost.

\begin{figure}[H]\centering
\figtitle{Accommodative monetary offset}
\IfFileExists{output/fig_monetary_offset_import.pdf}{\includegraphics[width=\textwidth]{output/fig_monetary_offset_import.pdf}}{\textit{[figure output/fig\_monetary\_offset\_import.pdf: run run\_monetary\_offset.py]}}
\label{fig:monetary_offset}
\fignotes{Run alongside the brown-price shock, sized to halve the cumulative output
loss: energy shock alone (dashed) against energy shock plus accommodation (solid).
Level deviations $\times100$.}
\end{figure}
